\documentclass[11pt,a4paper]{article}
\setlength{\oddsidemargin}{-0.44cm}
\setlength{\evensidemargin}{-0.44cm}
\setlength{\topmargin}{-0.44cm}
\setlength{\headheight}{0pt}
\setlength{\headsep}{0pt}
\setlength{\textwidth}{16.8cm}
\setlength{\textheight}{25cm}
\usepackage{amsmath,amssymb,graphicx,url}
\newcommand{\studentnumber}{STUDENTNUMBER}
\newsavebox{\resultbox}
\newcommand{\resulttable}[2]{%
  \par\medskip\noindent\begin{minipage}{\linewidth}\centering
  \textbf{Table #1. #2}\par\medskip\small
  \sbox{\resultbox}{\input{outputs/tex/Table_#1.tex}}%
  \ifdim\wd\resultbox>\linewidth
    \resizebox{\linewidth}{!}{\usebox{\resultbox}}%
  \else\usebox{\resultbox}\fi
  \end{minipage}\par\medskip}
\newcommand{\resultfigure}[2]{%
  \par\medskip\noindent\begin{minipage}{\linewidth}\centering
  \includegraphics[width=0.90\textwidth]{outputs/figures/Figure_#1.pdf}
  \par\small #2\end{minipage}\par\medskip}
\title{Assignment 2: Position-Aware Bandits for Fashion Recommendations}
\author{Student number: \studentnumber}
\date{Learning from Big Data 2026}
\begin{document}
\maketitle
\noindent\textbf{Working draft.} Replace the student number and review the interpretation before submission.
All numerical tables and figures below are generated by the accompanying Python program.
Rates in tables are fractions, so 0.005 means 0.5\% CTR. Figures label percentages explicitly.

\section{Data, estimand and experimental design}
The assignment asks for two extensions of TS/UCB and evaluation of CTR, batch size,
hyperparameters, heterogeneity, fairness and stationarity. The supplied file contains
logged item recommendations with binary clicks, positions, propensities and four hashed
categorical user features. A row is one displayed slot, not an entire three-item page view.
We maximize the sum of expected clicks over three distinct items assigned to three positions.
We assume a slot's reward does not depend on the other items in the slate, following the
dataset description~\cite{saito}. Under this assumption, marginal item-at-position OPE
can evaluate valid slate policies without reconstructing full impressions.

\paragraph{Retained action support.}
The file named ``80items'' contains 79 observed IDs (1--79), although all supplied
propensities are $1/80$. We restrict policies to observed items. Our explicit assumption
is that the file retained these items from a uniform 80-item log through item-only
filtering, without additional reward/context-dependent selection. The conditional
behavior propensity is then $b(a\mid x,p)=1/79$. This assumption is not verified by
the CSV alone. Table 3's full CSV also reports IPS with the original propensity; that
estimate differs by a factor $80/79$. SNIPS is unchanged by this common scaling.
None of these estimates establish performance on the original full 80-item population.
The original file is never modified.

There are no missing required values. We retain 162 duplicate-looking rows because
the file lacks an impression identifier that would establish whether these are errors.
Hashed user features are treated as categories, without attributing real demographics
to hash values. No user identifiers are available for clustering repeated visitors.

\resulttable{02}{Chronological split and logged CTR}
The first three days fit models, the next two days select hyperparameters, and the final
two days evaluate frozen policies. After selection, we refit on the first five days.
Test outcomes never select hyperparameters. A separate validation replay experiment
studies adaptive learning and delayed batch feedback. Its traffic clock and policy
estimand differ from the frozen-policy test experiment and are labelled accordingly.

\section{(a) Methods and comparison with the logging policy}
\subsection{Position-aware Thompson Sampling}
For each item $a$ and position $p$, let $N_{a,p}$ and $S_{a,p}$ denote displays and clicks
in the historical fitting sample. Define the empirical prior center
$m=(S_{\rm all}+1)/(N_{\rm all}+2)$ and prior strength $\kappa$. We sample
\[
\theta_{a,p}\sim\operatorname{Beta}\{S_{a,p}+\kappa m,
N_{a,p}-S_{a,p}+\kappa(1-m)\}.
\]
We then find the maximum-score assignment of three distinct items to the three slots.
Unlike selecting a separate best item independently for each slot, this construction
prevents duplicate recommendations within a page.
This is an empirical-Bayes TS extension, with a low-CTR prior estimated from fitting
data rather than an uninformative 50\% prior mean.

In the test experiment, each of five fixed seeds produces 256 posterior-sampled slates.
The evaluation policy is defined as their uniform finite mixture, plus the exploration
branch below. Its marginal action probabilities are therefore known exactly for the
constructed finite policy; it approximates the ideal integrated TS distribution.
The posterior remains frozen during test.

\subsection{Position-aware smoothed UCB}
We use the same item--position statistics and optimize the assignment using
\[
U_{a,p}=\frac{S_{a,p}+\kappa m}{N_{a,p}+\kappa}
+c\sqrt{\frac{\log(2+\sum_{a,p}N_{a,p})}{N_{a,p}+1}}.
\]
The smoothed mean handles sparse click observations, while $c$ controls exploration.
This is a UCB-derived heuristic; we do not claim the exact regret guarantee of classical
UCB1 for this modified score or its finite-log evaluation.

For both algorithms, with probability $\epsilon=0.05$ we select a uniformly random
distinct-item slate; otherwise we use the algorithm's slate. The resulting marginal
probability at each slot is
$\pi=(1-\epsilon)\pi_{\rm base}+\epsilon/K$, $K=79$.
The implementation solves the assignment with \texttt{linear\_sum\_assignment}.

\subsection{Estimators and uncertainty}
We observe a reward only for the logged action. With
$w_i=\pi(a_i\mid x_i,p_i)/b(a_i\mid x_i,p_i)$, the estimators are
\begin{align*}
\widehat V_{\rm IPS}&=\frac1n\sum_i w_i r_i,&
\widehat V_{\rm SNIPS}&=\frac{\sum_i w_i r_i}{\sum_i w_i},\\
\widehat V_{\rm DR}&=\frac1n\sum_i\left[
\sum_a\pi(a\mid x_i,p_i)\hat q(a,p_i)
+w_i\{r_i-\hat q(a_i,p_i)\}\right].
\end{align*}
The DR reward model is a smoothed historical item--position mean fitted before the
evaluation period. DR is the primary measure, accompanied by IPS, SNIPS and effective
sample size $\operatorname{ESS}=(\sum w_i)^2/\sum w_i^2$. We do not clip weights or
negative estimated values. The validity of OPE still depends on logging propensities,
support and the stated sampling assumptions~\cite{saito}.

The reference ``Logger'' is the observed CTR on the same test rows, consistent with
uniform logging in the retained item set. ``Pooled TS'' uses item-only statistics and is
a useful learning baseline, not an exact reproduction of Zozo's production Bernoulli TS.
Production TS results from different 10/20-item files are not a comparable control.

Uncertainty uses 2,000 paired resamples of non-overlapping hourly blocks of test DR
contributions. The same sampled blocks are used for both policies in every contrast.
These are approximate intervals conditional on the fitted policies and reward model;
training uncertainty, long-range visitor dependence and repeated-testing corrections
are not covered. The two-day test period limits long-term generalization.

\resulttable{03}{Frozen policy values and conditional uncertainty on test}
\resulttable{04}{Paired DR contrasts on identical test rows}
\resultfigure{01}{Intervals are 95\% hourly-block bootstrap intervals. They are not variation across random seeds.}

\subsection{Automatically generated findings}
The supplied archive contains 1,356,670 slot observations and 79 observed items. The original propensity is 0.0125. Results condition on the retained item universe under the item-only filtering assumption.

Logger: estimated DR CTR 0.361\% (95\% hourly-block interval [0.329, 0.392]\%), ESS 378,533.

Pooled TS: estimated DR CTR 0.677\% (95\% hourly-block interval [0.541, 0.818]\%), ESS 24,799.

Position TS: estimated DR CTR 0.585\% (95\% hourly-block interval [0.441, 0.739]\%), ESS 10,700.

Position UCB: estimated DR CTR 0.488\% (95\% hourly-block interval [0.303, 0.689]\%), ESS 5,622.

Fair TS: estimated DR CTR 0.504\% (95\% hourly-block interval [0.398, 0.611]\%), ESS 37,047.

Pooled TS - Logger: DR difference 0.317 percentage points; 95\% interval [0.191, 0.450] percentage points, positive.

Position TS - Logger: DR difference 0.224 percentage points; 95\% interval [0.078, 0.383] percentage points, positive.

Position UCB - Logger: DR difference 0.128 percentage points; 95\% interval [-0.068, 0.338] percentage points, inconclusive in sign.

Fair TS - Logger: DR difference 0.144 percentage points; 95\% interval [0.051, 0.239] percentage points, positive.

Position TS - Position UCB: DR difference 0.097 percentage points; 95\% interval [-0.026, 0.207] percentage points, inconclusive in sign.

The fairness extension changes item exposure Gini from 0.870 to 0.419. Worst supported group uplift changes from -107.9\% to -77.1\%. These are descriptive group point estimates, not a statistical or ethical guarantee of user fairness.

Unadjusted reward by day: adjusted q = 0.000999; rejects the tested stability null at 5\%. Non-rejection does not establish stationarity.

Reward given item and position: adjusted q = 0.000999; rejects the tested stability null at 5\%. Non-rejection does not establish stationarity.

Reward given item position and feature 0: adjusted q = 0.000999; rejects the tested stability null at 5\%. Non-rejection does not establish stationarity.

Context mix by day: user\_feature\_0: adjusted q = 1.14e-231; rejects the tested stability null at 5\%. Non-rejection does not establish stationarity.

Context mix by day: user\_feature\_1: adjusted q = 1.851e-183; rejects the tested stability null at 5\%. Non-rejection does not establish stationarity.

Context mix by day: user\_feature\_2: adjusted q = 4.11e-172; rejects the tested stability null at 5\%. Non-rejection does not establish stationarity.

Context mix by day: user\_feature\_3: adjusted q = 3.271e-278; rejects the tested stability null at 5\%. Non-rejection does not establish stationarity.


The point estimate alone must not determine a winner when the paired difference interval
contains zero. The item-only TS baseline can outperform a more detailed method when
additional position or group parameters create variance with sparse rewards. Improvements
in modeling detail are hypotheses to test, not guaranteed improvements in recommendation quality.

\section{(b) Sensitivity to batch size}
We warm-start each algorithm from the first three days, then process validation logs
chronologically. Within each batch, the policy is fixed before observing batch rewards.
Only matched feedback updates the policy: a logged row is accepted with probability
$\pi(a_i\mid x_i,p_i)$. This is equivalent to sampling a target action and checking if it
matches the logged action. Under uniform logging, marginal acceptance is $1/K$, which
preserves the context distribution in accepted events~\cite{li}. The posterior/counts
are updated at the end of the batch, including any final partial batch.

Batch sizes are 500, 5,000 and 50,000 \emph{logged slot rows}; they are not page views,
accepted feedback events or time in seconds. A 5,000-row batch yields approximately
$5000/79$ accepted events. The experiment represents delayed replay on the supplied log
clock, not a full reconstruction of deployed online learning speed. We report replay CTR
and pre-update cumulative IPS; the latter uses probabilities determined before rewards.
Five seeds describe algorithmic randomness. Seed standard deviations are not confidence
intervals for deployment CTR.

\resulttable{05}{Replay batch-size sensitivity on validation}
\resultfigure{02}{Error bars show standard deviation across seeds, not 95\% confidence intervals.}
\resultfigure{07}{One-seed illustration at 5,000 logged rows per update; it does not establish statistical superiority.}
Small batches permit more frequent adaptation, while larger batches retain stale scores
longer. With rare clicks and only a small fraction of matched feedback, the measured
relationship need not be monotonic. The observed means and seed spread should be discussed
together; an apparently best batch on validation is not independently confirmed on test.

\section{(c) Sensitivity to parameter tuning}
For TS we compare $\kappa\in\{10,100,1000\}$. For UCB we compare
$c\in\{0.001,0.01,0.1,1\}$ with $\kappa=100$. Selection maximizes mean validation DR
over five seeds. The exploration mixture remains fixed at 0.05 in these experiments.
Large $\kappa$ increases shrinkage; large $c$ emphasizes under-observed pairs.
The selected value is best only within this predefined grid and validation period.

\resulttable{06}{Parameter sensitivity on validation}
\resultfigure{03}{TS prior strength and UCB exploration coefficient use different axes and retain their own meanings.}
\resulttable{14}{Final configurations locked before test evaluation}

\section{(d) Aggregation and heterogeneity}
We compare item-only (pooled), item--position, and item--position--group models.
The group experiment uses user feature 0. All alternatives use the same validation
period and the family-specific selected hyperparameter. For a group model, the prior
center for a cell is the global smoothed item--position estimate fitted on training data.
Unknown groups fall back to the global model.

This partial pooling is an empirical-Bayes approximation, not exact hierarchical
Bayesian inference: the parent estimates include group observations. It stabilizes sparse
cells but adds a modeling assumption. Keeping a single grouping variable also avoids
an extremely sparse full Cartesian product of all four features.

\resulttable{07}{Aggregation and heterogeneity comparison}
\resultfigure{04}{Extra heterogeneity can trade lower approximation bias for higher estimation variance.}
The analysis should report whether the finer model actually helps, without assuming
that personalization must improve results. Differences between seed values can also
arise from breaking assignment ties when a pooled score cannot distinguish positions.

\section{(e) Fairness and an implemented extension}
\subsection{Fairness to users}
We audit every level of each of the four hashed features separately, as permitted by
the assignment. Measures include group DR/SNIPS CTR, the largest within-feature CTR
gap, and worst supported group uplift relative to that group's own logged CTR.
The latter distinguishes utility changes from differences in baseline click propensity.
Equal observed CTR across groups is not automatically an appropriate normative objective.

All groups are retained in Table 9's CSV. Summary comparisons use the same eligible
groups for every policy: at least $100K=7900$ logged rows and 10 logged clicks.
Realized ESS below 100 is flagged separately. These fixed thresholds do not eliminate
estimation noise, and feature intersections are not fully audited.
Group results are descriptive point estimates, not simultaneous confidence bounds.

In particular, DR can be negative in a sparse subgroup, producing an estimated relative
uplift below $-100\%$. This is sampling/model noise in an unbounded estimator, not a
possible true negative click probability. Such values must be read alongside bounded
SNIPS estimates and ESS in Table 9. They prevent a strong claim of universal user benefit.

\subsection{Fairness to items}
Expected item exposure is computed from policy probabilities on all evaluation contexts,
not from the subset of matched replay rows. We include all 79 items, even if their share
is zero. We report exposure Gini, normalized entropy, minimum/maximum share and coverage.
This is equal opportunity for exposure, not a merit-adjusted or seller-level fairness
definition. With nonzero uniform exploration, coverage is always 100\%, so Gini and
minimum share provide more discrimination.

\subsection{Fair TS}
We implement
\[
\pi_{\rm fair}(a\mid x,p)=(1-\epsilon)\pi_{\rm segmented\ TS}(a\mid x,p)+\epsilon/K.
\]
Validation tries each of the four features and $\epsilon\in\{0.10,0.25,0.50\}$.
Among candidates retaining at least 95\% of Position TS's validation DR, select the
largest worst supported group DR uplift. Table 16 records the full selection audit.
If no candidate is feasible, code selects a labelled best-effort alternative and records
that fact in the manifest. The empirical utility constraint is not a population guarantee.

The uniform branch guarantees marginal exposure of at least $\epsilon/K$ per item per
slot. Segmentation seeks to improve utility across groups, but does not mathematically
guarantee user parity, no harm to every group, or the same utility floor on test.
Assess actual test outcomes and distinguish the item-exposure guarantee from user
utility estimates. Gini also changes because personalization itself changes allocations;
the entire difference should not be attributed to epsilon alone.

\resulttable{08}{Fairness summary on the common test population}
\resultfigure{05}{Group uplift values are noisy point estimates; an estimate below $-100\%$ is not a feasible true utility change.}
The extension trades some aggregate CTR for a more even allocation of item exposure.
Whether that trade-off is desirable depends on the chosen fairness objective. The present
results do not certify demographic fairness or improvement for every protected group.

\section{(f) Empirical stationarity evidence}
\resulttable{10}{Daily logged CTR and descriptive Wilson intervals}
\resultfigure{06}{Day-to-day CTR is descriptive; conditional tests help separate reward changes from changing exposure/context mix.}
We test constant click probabilities across days using a binomial likelihood-ratio
statistic, first unadjusted, then within item--position cells, and then within
item--position--feature-0 cells. The alternative allows separate daily probabilities.
For sparse cells, calibration uses 1,000 parametric binomial bootstrap samples under
the pooled null, holding cell/day impression counts fixed and re-estimating the null
for each bootstrap sample. The reported p-value is $(1+\#\{G^*\ge G\})/(1000+1)$.
Thus the simulation resolution is approximately 0.001.

We also test each user-feature distribution against day with a contingency-table
chi-square test; tests with very sparse category counts are exploratory. We apply
Benjamini--Hochberg adjustment across these seven tests. The tests assume conditionally
independent Bernoulli/categorical observations. Repeated users and residual serial
dependence cannot be fully controlled, and may make significance appear stronger.

\resulttable{11}{Reward and context stability tests}
Reward drift and context drift are different phenomena. Conditional reward tests are
more informative than the raw CTR curve about changes beyond the displayed-item mix,
but do not adjust for every unobserved factor. Rejection is evidence against the tested
stability null under its assumptions; non-rejection would not prove stationarity.
Seven days of evidence does not establish the behavior over seasons or future campaigns.

\section{(g) Model changes for non-stationarity}
An implemented example exponentially discounts historical counts at each batch:
\[
N\leftarrow\gamma N+N_{\rm batch},\qquad
S\leftarrow\gamma S+S_{\rm batch}.
\]
The prior remains fixed. This discounts warm-start data as well as older replay feedback.
We compare $\gamma\in\{1,0.99,0.95\}$ on validation, always updating after 5,000 logged
rows. For $0<\gamma<1$, the count half-life is $\log(0.5)/\log(\gamma)$ batches.
The table is a sensitivity illustration, not a jointly tuned gamma/batch deployment test.

\resulttable{15}{Discounted TS/UCB replay sensitivity}
Other changes include a sliding time window, residual-based change detection followed
by a reset or exploration increase, and time-of-day/day-of-week context. Each introduces
choices about window length, detection threshold or context complexity that must be
selected on past-to-future validation splits. Smaller memory adapts faster but increases
variance for low-CTR items. Drift evidence alone does not establish that discounting is
better than the stationary model; the empirical comparison is necessary.

\section{Reproducibility and limitations}
Run \texttt{python run\_assignment.py} from this directory, or use the repository-relative
path in the guide. The default uses every source row, fixes seed 1 and evaluates seeds
1--5. Every table is printed with its identifier and exported to CSV; every figure is
exported as PNG/PDF with a printed identifier and path. Use \texttt{--show} to open figures.
\texttt{--quick} is explicitly a subsampled demonstration and must not supply the final report.

The manifest records the source checksum, versions, settings, support assumptions and
fairness-selection feasibility. Tables are included directly from generated LaTeX snippets
to avoid transcription errors. Tables 1, 9, 12, 13, 16 and 17 provide full audits in CSV
and console output; Table 12 provides illustrative posterior-mean item/position assignments,
which differ from the stochastic policy evaluated in Table 3.

The main limitations are unverified item filtering, sparse click outcomes, short evaluation
time, unavailable visitor/impression identifiers, group-estimate instability, possible
non-stationarity, and the no-interference assumption across slots. Counterfactual policy
estimates and replay are not a substitute for a future randomized online experiment.

\noindent\begin{minipage}{\linewidth}
\begin{thebibliography}{9}
\bibitem{saito} Y. Saito, S. Aihara, M. Matsutani and Y. Narita (2021).
\emph{Open Bandit Dataset and Pipeline: Towards Realistic and Reproducible Off-Policy Evaluation}.
NeurIPS Datasets and Benchmarks. Supplied as \texttt{zozo paper.pdf}.
\url{https://arxiv.org/abs/2008.07146}.
\bibitem{li} L. Li, W. Chu, J. Langford and X. Wang (2011).
\emph{Unbiased Offline Evaluation of Contextual-bandit-based News Article Recommendation Algorithms}.
WSDM. Preprint posted in 2010. \url{https://arxiv.org/abs/1003.5956}.
\bibitem{obp} ZOZO Research. Open Bandit Pipeline source and dataset documentation.
\url{https://github.com/st-tech/zr-obp}.
\end{thebibliography}
\end{minipage}
\end{document}
